\section*{第 \ref{ch:b2:dipole-radiation} 章——偶极辐射与散射}

\begin{solution}{exo:b2:dipole-radiation:1}
$p_0 = I_0\ell/\omega = \qty{3.2e-9}{C.m}$; $E = \mu_0\omega^2p_0\sin\theta/4\pi r$: \qty{0.13}{V/m}
在 $\theta = 90^\circ$ 时如上，\qty{30}{\degree} 时为 \qty{0.063}{V/m}；$\mathcal P = \mu_0\omega^4p_0^2/12\pi c =
\qty{175}{W}$；$R_{\text{辐射}} = 2\mathcal P/I_0^2 = \qty{88}{\ohm}$。
\end{solution}

\begin{solution}{exo:b2:dipole-radiation:2}
$790(\ell/\lambda)^2$：\qty{8.8e-5}{\ohm}、\qty{8.8e-3}{\ohm}、\qty{0.88}{\ohm}、\qty{88}{\ohm}。铜：
$0.026$、$0.084$、$0.27$、\qty{0.84}{\ohm}（趋肤深度）：效率分别为 $0.3\%$、$10\%$、
$77\%$、$99\%$。波长很长时，只有高度占
$\lambda$ 相当大比例的桅杆，其辐射电阻才超过其损耗。
\end{solution}

\begin{solution}{exo:b2:dipole-radiation:3}
$400 : 550 : 700 = 9.4 : 2.6 : 1$。$(10/3)^4 = 120$：\qty{3}{cm} 雷达
探测小水滴的效果好百倍；\qty{10}{cm} 雷达
则可穿过雨区，看到后方的目标。
\end{solution}

\begin{solution}{exo:b2:dipole-radiation:4}
$\sin^2\chi/(1 + \cos^2\chi)$：$0.14$、$0.60$、$1$、$0.60$。观测点与太阳成
\qty{90}{\degree} 角（天空中的一个大圆）；太阳位于头顶时，整个
地平线处于 \qty{90}{\degree} 方向，并沿之偏振。多次散射、
地面反射与分子各向异性共同贡献非偏振光，
至多约 $75\%$。
\end{solution}

\begin{solution}{exo:b2:dipole-radiation:5}
(a) $\int\sin^2\theta\,\dd\Omega = 8\pi/3$. (b) $\int_{60^\circ}^{120^\circ}\sin^3\theta\,\dd\theta/
(4/3) = 0.917/1.333 = 69\%$。(c) $\sin^2\theta = \tfrac12$：$\theta = \qty{45}{\degree}$，一束
光束宽 \qty{90}{\degree}。(d) 最大值 $1$ 与平均值 $2/3$ 之比为 $1.5$。
\end{solution}

\begin{solution}{exo:b2:dipole-radiation:6}
(a) $p/4\pi\varepsilon_0r^3 = \omega^2p/4\pi\varepsilon_0c^2r$，在 $r_0 = c/\omega = \lambda/2\pi$ 处。(b) \qty{48}{m}：
在 \qty{10}{m} 处为准静态场，在 \qty{10}{km} 处为辐射场。(c)
$r_0 = \qty{4.8}{cm}$：头部位于近区。(d) 近区中的 $\vect E$
与 $\vect B$ 正交：坡印廷矢量的平均值为零，
能量储存后又被返还，与电容器中的情形相同。
\end{solution}

\begin{solution}{exo:b2:dipole-radiation:7}
(a) $8\pi r_e^2/3 = \qty{6.65e-29}{m^2}$. (b) $n_e\sigma_TL = 10^{14} \times 6.65 \times 10^{-29}
\times 10^9 = 7 \times 10^{-6}$：约为光球层的百万分之一，因汤姆孙散射本身
无色，且被天空散射的阳光所淹没，
除日食时外。(c) \qty{10}{keV} $\gg$ \qty{300}{eV}：电子的响应如同
自由电子，$6\sigma_T = \qty{4e-28}{m^2}$。(d) 光电吸收按 $Z^4$ 增长，
并使钙尤为突出；散射对骨骼与
软组织中每个电子的作用相同。
\end{solution}

\begin{solution}{exo:b2:dipole-radiation:8}
(a) $\tau = 0.36$, $0.10$, $0.038$: $T = 0.70$, $0.90$, $0.96$. (b) $\tau \times 38 = 13.6$,
$3.8$、$1.45$：$T = 1 \times 10^{-6}$、$0.02$、$0.23$；红光与蓝光之比 $\sim 2 \times 10^5$。(c)
地球大气将日落时的光线弯入阴影区，
并由同样的散射而变红。(d) 近地平线处光程很长：蓝光
从散射光本身中被散射掉，多次散射又使
光变白。
\end{solution}

\begin{solution}{exo:b2:dipole-radiation:9}
(a) 沿这对天线的连线方向，光程差 $\lambda/2$ 导致
相消；侧向则相加。(b) 反相时情形相反——
沿连线方向（端射）辐射，侧向无辐射。(c) 朝 $+x$ 方向：
光程延迟 $-\pi/2$ 加上馈电相位 $+\pi/2$，同相；朝 $-x$ 方向：
$-\pi$，相消：呈心形，单侧辐射。(d) 天线桅杆决定覆盖范围
并保护邻近设施；相控阵通过改变馈电
相位来转向，只需微秒量级。
\end{solution}

\begin{solution}{exo:b2:dipole-radiation:10}
(a) $\langle\mathcal P\rangle = e^2\omega_0^4x_0^2/12\pi\varepsilon_0c^3$, $W = \tfrac12m\omega_0^2x_0^2$:
$\Gamma = \mathcal P/W = e^2\omega_0^2/6\pi\varepsilon_0mc^3$. (b) $\omega_0 = \qty{3.2e15}{rad/s}$: $\Gamma =
\qty{6.4e7}{s^{-1}}$——即所用的值。(c) \qty{16}{ns}；$\Delta\lambda = \lambda^2\Gamma/2\pi c =
\qty{1.2e-14}{m}$，约为百分之一皮米（\qty{10}{MHz}）。(d) $Q = 5 \times 10^7$。
\end{solution}

\begin{solution}{exo:b2:dipole-radiation:11}
(a) \qty{75}{m}; $I_0 = \sqrt{2\mathcal P/R} = \qty{53}{A}$. (b) $\langle\Pi\rangle = 1.5 \times
5 \times 10^4/2\pi \times 10^8 = \qty{1.2e-4}{W/m^2}$, $E_0 = \sqrt{2\Pi/\varepsilon_0c} = \qty{0.30}{V/m}$.
(c) \qty{0.30}{V}。(d) 回流电流在大地中流动；电阻性的
土壤会增加损耗、破坏镜像作用；埋设的
地网为其提供铜质通路。
\end{solution}

\begin{solution}{exo:b2:dipole-radiation:12}
(a) 由 $p_0 = \alpha E_0$ 得 $\sigma = \alpha^2\omega^4/6\pi\varepsilon_0^2c^4$；其中 $\alpha = \varepsilon_0(n^2
- 1)/N$，$\omega = 2\pi c/\lambda$，于是 $\sigma = 8\pi^3(n^2 - 1)^2/3N^2\lambda^4$。(b) $n^2 - 1 =
5.8 \times 10^{-4}$：$\sigma = \qty{4.9e-31}{m^2}$，平均自由程 $1/N\sigma = \qty{80}{km}$。(c)
$8/80 = 0.1$。(d) $\sigma$ 由天空的衰减测得，$n$ 已知，
$N$ 随之确定，进而得阿伏伽德罗常数。
\end{solution}

\begin{solution}{pb:b2:dipole-radiation:1}
\textbf{1.} $\lambda = \qty{300}{m}$, $h = \qty{75}{m}$; $I_0 = \sqrt{2 \times 10^5/36} = \qty{75}{A}$;
$V = RI_0 = \qty{2.7}{kV}$.

\textbf{2.} $\ell_{\text{有效}} = \qty{48}{m}$：$p_0 = 75 \times 48/6.3 \times 10^6 = \qty{5.7e-4}{C.m}$。

\textbf{3.} $\langle\Pi\rangle = 1.5 \times 10^5/2\pi r^2$：\qty{2.4e-4}{W/m^2} 与 \qty{2.4e-6}{W/m^2}；
$E_0 = \qty{0.42}{V/m}$ 与 \qty{0.042}{V/m}。

\textbf{4.} 鞭状天线：\qty{42}{mV}。磁棒天线：$N\mu_{\text{有效}}A\omega B_0 = 10^4 \times 10^{-4} \times 6.3 \times
10^6 \times 1.4 \times 10^{-10} = \qty{0.9}{mV}$——更小，但更紧凑、更具选择性。

\textbf{5.} $A = 1.5\lambda^2/4\pi = \qty{1.1e4}{m^2}$: $P = \Pi A = \qty{26}{mW}$.

\textbf{6.} $36/38 = 95\%$；\qty{5}{kW} 的热量。

\textbf{7.} 在 \qty{10}{km} 处 $r/\lambda = 33$，为辐射区；在 \qty{100}{m} 处，
$r \approx \lambda/3$，为过渡区，此处准静态场仍
起作用。

\textbf{8.} 白天，D 层吸收本会
抵达反射层的波；夜间 D 层消失，天波
可将电台信号传送数百千米。

\textbf{9.} $\sigma(450) = \qty{1.0e-30}{m^2}$，$\sigma(650) = \qty{2.4e-31}{m^2}$；比值为 $4.4$。

\textbf{10.} 平均自由程分别为 $39$、$87$、\qty{170}{km}；$\tau = 0.21$、$0.09$、$0.05$：
$19\%$、$9\%$、$5\%$ 被散射。

\textbf{11.} $N\sigma I = 2.5 \times 10^{25} \times 4.6 \times 10^{-31} \times 300 = \qty{3.5}{mW/m^3}$，
对自然光散射至 $4\pi$ 立体角，方向图为 $(1 + \cos^2\chi)$。

\textbf{12.} 该光柱散射 $3.5 \times 10^{-3} \times 8000 = \qty{28}{W}$，占其
\qty{300}{W} 的 $9\%$；若均匀分布于半球，天空提供约
\qty{4}{W.m^{-2}.sr^{-1}}，而太阳为 $300/6.8 \times 10^{-5} = \qty{4e6}{W.m^{-2}.sr^{-1}}$：
即每球面度仅为太阳的百万分之一，量级正确。

\textbf{13.} 单次散射时，在 \qty{90}{\degree} 处完全偏振；多次
散射及地面反射使其降至约 $75\%$；蜜蜂和蚂蚁
读取偏振图样，即使云层遮挡也能找到太阳。

\textbf{14.} $\tau \times 38$：蓝光 $8$（$T = 3 \times 10^{-4}$），红光 $1.8$（$T = 0.17$）：太阳呈深
红色；头顶上方的光路短，呈蓝色。

\textbf{15.} 无大气，无物可散射：太阳
旁的天空为黑色。

\textbf{16.} 微米级尘埃像米氏粒子那样散射：前向
散射偏红，使天空呈奶油糖色；在落日附近，
前向散射光是蓝色的。

\textbf{17.} $8\pi r_e^2/3 = \qty{6.65e-29}{m^2}$。

\textbf{18.} $a = v^2/2d = \qty{6.4e21}{m/s^2}$；$\mathcal P = e^2a^2/6\pi\varepsilon_0c^3 =
\qty{2e-10}{W}$，持续 $2d/v = \qty{2.5e-14}{s}$，辐射 \qty{6e-24}{J}，
为 \qty{100}{keV} 的 $10^{-10}$；实际产额约为 $1\%$，
这是因为制动发生在与原子核的剧烈近距碰撞中，
而非平缓的匀减速。

\textbf{19.} $hc/E = \qty{12.4}{pm}$。

\textbf{20.} 束流功率 \qty{2}{kW}，其中 X 射线约 \qty{20}{W}；
焦斑上的热量同样是两千瓦——阳极旋转以分散热量。

\textbf{21.} $a = v^2/r = \qty{9e22}{m/s^2}$，$\mathcal P = \qty{5e-8}{W}$；\qty{13.6}{eV} 的能量
在 \qty{5e-11}{s} 内耗尽。

\textbf{22.} 束缚运动中的经典电荷持续辐射；
量子力学中的定态不辐射，并有一个最低能级，
无处可以跃迁下去。

\textbf{23.} $a = c^2/R = \qty{9e14}{m/s^2}$：由拉莫尔公式得 \qty{5e-24}{W}，实际上多出 $\gamma^4 \sim 10^{16}$
倍：存储束流所产生的同步辐射光可达数十千瓦。
%

\textbf{24.} $\sigma_TI = \qty{6.7e-26}{W}$；对应一瓦的电子数为 $1.5 \times 10^{25}$。

\textbf{25.} 天线杆：\qty{5.7e-4}{C.m}，频率 \qty{1}{MHz}，\qty{100}{kW}，竖直
环形方向图。分子：$\sim$\qty{1e-35}{C.m}，频率 \qty{500}{THz}，\qty{1e-28}{W}，
围绕场的环形方向图。自由电子：汤姆孙散射，
日光下 \qty{1e-25}{W}，方向图相同。
\end{solution}
